\chapter{Boolean Analysis — Feed Forward Circuit}
%%% generated by the blueprint pipeline from TCSlib.BooleanAnalysis.RazborovSmolensky.FeedForwardCircuit
%%%%%%%%%%%%%%%%%%%%%%%%%%%%%%%%%%%%%%%%%%%%%%%%%%%%%%%%%%%%%%%%%%%%%%%%%%%%%%%
\section{Overview}

This file introduces layered feedforward circuits over an arbitrary alphabet $\alpha$:
gates with explicit input wiring, the evaluation of nodes and of the whole circuit,
and the basic complexity measures (depth, size, finiteness, gate set restriction).
It then relates feedforward circuits to the tree-shaped \texttt{BoolCircuit.Circuit}
in both directions: tree-unrolling a layered AND/OR circuit into a tree (with an
exponential-in-depth size bound), and embedding a tree circuit back as a
feedforward circuit padded with identity wires.

%%%%%%%%%%%%%%%%%%%%%%%%%%%%%%%%%%%%%%%%%%%%%%%%%%%%%%%%%%%%%%%%%%%%%%%%%%%%%%%
\section{Declarations}

\begin{definition}[Gate with input wiring]
\lean{ACP.Gate}
\label{ACP.Gate}
\leanok
A gate over alphabet $\alpha$ with input domain \texttt{domain} consists of a gate
operation \texttt{op} (an arity type $\iota$ together with a function
$(\iota \to \alpha) \to \alpha$) and a wiring map $\iota \to \texttt{domain}$ that
assigns to each input slot the node it reads from.
\end{definition}

\begin{definition}[Layered feedforward circuit]
\lean{ACP.FeedForward}
\label{ACP.FeedForward}
\leanok
\uses{ACP.Gate}
A feedforward circuit from \texttt{inp} to \texttt{out} over $\alpha$ consists of a
depth $d \in \bbn$, a family of node types indexed by the layers
$\mathrm{Fin}\,(d+1)$, and, for every layer $i < d$, a gate for each node of layer
$i+1$ whose inputs are wired to nodes of layer $i$. Layer $0$ is required to be the
input type \texttt{inp} and the last layer the output type \texttt{out}.
\end{definition}

\begin{definition}[Identity gate operation]
\lean{ACP.FeedForward.GateOp.id}
\label{ACP.FeedForward.GateOp.id}
\leanok
The gate operation of arity $\mathrm{PUnit}$ that returns its single input
unchanged.
\end{definition}

\begin{definition}[Evaluation of a single gate]
\lean{ACP.FeedForward.Gate.eval}
\label{ACP.FeedForward.Gate.eval}
\leanok
\uses{ACP.Gate}
Given a gate $g$ and an assignment $xs$ of values to the nodes of its input domain,
the value of $g$ is $g.\texttt{op}.\texttt{func}\,(xs \circ g.\texttt{inputs})$, i.e.\
the gate operation applied to the values read along the input wires.
\end{definition}

\begin{definition}[Evaluation of a node]
\lean{ACP.FeedForward.evalNode}
\label{ACP.FeedForward.evalNode}
\leanok
\uses{ACP.FeedForward.Gate.eval}
For an input assignment $xs : \texttt{inp} \to \alpha$, the value of a node at layer
$d$ is defined by recursion on $d$: a layer-$0$ node is read off directly from $xs$,
and a node at layer $n+1$ evaluates its gate on the values of the layer-$n$ nodes.
\end{definition}

\begin{definition}[Evaluation of a circuit]
\lean{ACP.FeedForward.eval}
\label{ACP.FeedForward.eval}
\leanok
\uses{ACP.FeedForward.evalNode}
The function $\texttt{out} \to \alpha$ sending each output node to its value in the
last layer under the input assignment $xs$.
\end{definition}

\begin{definition}[Evaluation with a unique output]
\lean{ACP.FeedForward.eval₁}
\label{ACP.FeedForward.eval₁}
\leanok
\uses{ACP.FeedForward.eval}
When the output type is a \texttt{Unique} type, the single value
$F.\texttt{eval}\,xs\,\texttt{default} \in \alpha$ computed by the circuit.
\end{definition}

\begin{definition}[Size of a feedforward circuit]
\lean{ACP.FeedForward.size}
\label{ACP.FeedForward.size}
\leanok
The size of $F$ is the number of its non-input nodes, i.e.\ the total number of
nodes lying in layers $1, \dots, \mathrm{depth}(F)$.
\end{definition}

\begin{definition}[Finiteness of all layers]
\lean{ACP.FeedForward.Finite}
\label{ACP.FeedForward.Finite}
\leanok
The proposition that the node type of every layer $i$ is finite.
\end{definition}

\begin{definition}[Restriction to a gate set]
\lean{ACP.FeedForward.onlyUsesGates}
\label{ACP.FeedForward.onlyUsesGates}
\leanok
For a set $S$ of gate operations, $\texttt{onlyUsesGates}\ S$ asserts that the
operation of every gate of the circuit belongs to $S$.
\end{definition}

\begin{theorem}[Every circuit has size at least one]
\lean{BoolCircuit.Circuit.one_le_size}
\label{BoolCircuit.Circuit.one_le_size}
\leanok
\uses{BoolCircuit.Circuit, BoolCircuit.Circuit.size}
\difficulty{2}
Let $C$ be a Boolean circuit on $n$ variables. Then the size of $C$ — the total number
of literal leaves and internal AND/OR gates — is at least $1$; that is, $1 \le
\mathrm{size}(C)$.
\end{theorem}

\begin{definition}[AND/OR gate restriction]
\lean{ACP.FeedForward.IsAndOrGate}
\label{ACP.FeedForward.IsAndOrGate}
\leanok
\uses{ACP.FeedForward}
Given a Boolean feedforward circuit $F$, a flag $\texttt{isAnd}$ on each non-input
node and a finiteness witness $\texttt{gfin}$ for each gate's arity,
$F.\texttt{IsAndOrGate}$ asserts that every gate operation is the conjunction of its
inputs when the flag is \texttt{true} and their disjunction when it is \texttt{false},
the inputs being enumerated by $\texttt{gfin}$. This is the gate restriction under
which $F$ can be converted into a \texttt{BoolCircuit.Circuit}.
\end{definition}

\begin{definition}[Tree-unrolling of a node]
\lean{ACP.nodeToCircuit}
\label{ACP.nodeToCircuit}
\leanok
\uses{ACP.FeedForward, BoolCircuit.Circuit}
Recursively expands a node $v$ of layer $m$ of a Boolean feedforward circuit into a
tree-shaped $\texttt{Circuit}\ n$: a layer-$0$ node becomes the positive literal of the
corresponding input variable, and a node at a higher layer becomes an AND/OR node
whose children are the unrollings of the nodes on its input wires. Nodes feeding
several downstream gates are duplicated.
\end{definition}

\begin{theorem}[Tree-unrolling preserves the value of a node]
\lean{ACP.nodeToCircuit_eval}
\label{ACP.nodeToCircuit_eval}
\leanok
\uses{ACP.FeedForward, ACP.FeedForward.Gate.eval, ACP.FeedForward.IsAndOrGate, ACP.FeedForward.evalNode, ACP.nodeToCircuit, BoolCircuit.Circuit.eval, BoolCircuit.Lit.eval}
\difficulty{5}
Let $F$ be a Boolean feedforward circuit whose inputs are indexed by $\mathrm{Fin}\,n$,
and suppose each of its gates has finite arity and computes either the conjunction or
the disjunction of its inputs, according to a flag marking every non-input node as an
AND node or an OR node. Then for every layer $m$ with $m \le \mathrm{depth}(F)$, every
node $v$ of layer $m$, and every input assignment $x : \mathrm{Fin}\,n \to \bbf_2$, the
Boolean circuit obtained by tree-unrolling $v$, evaluated at $x$, equals the value that
$v$ takes in $F$ under the assignment $x$.
\end{theorem}

\begin{theorem}[Exponential size bound for tree-unrolling]
\lean{ACP.nodeToCircuit_size_le}
\label{ACP.nodeToCircuit_size_le}
\leanok
\uses{ACP.FeedForward, ACP.nodeToCircuit, BoolCircuit.Circuit.size}
\difficulty{6}
Let $F$ be a Boolean feedforward circuit from the inputs $\mathrm{Fin}\,n$ to an output
type, in which every gate has finitely many input wires, and suppose $k \in \bbn$ bounds
these arities, so that each gate has at most $k$ input wires. Then for every layer $m
\le \mathrm{depth}(F)$ and every node $v$ of layer $m$, the tree-unrolling of $v$ is a
Boolean circuit of size at most $(k+1)^m$.
\end{theorem}

\begin{definition}[Conversion to a Boolean circuit]
\lean{ACP.FeedForward.toCircuit}
\label{ACP.FeedForward.toCircuit}
\leanok
\uses{ACP.FeedForward, ACP.nodeToCircuit, BoolCircuit.Circuit}
For an AND/OR feedforward circuit $F$ and a chosen output node $o : \texttt{out}$,
the tree-shaped $\texttt{Circuit}\ n$ obtained by tree-unrolling $o$ from the last
layer down to the inputs.
\end{definition}

\begin{theorem}[Tree-unrolling preserves evaluation]
\lean{ACP.FeedForward.toCircuit_eval}
\label{ACP.FeedForward.toCircuit_eval}
\leanok
\uses{ACP.FeedForward, ACP.FeedForward.IsAndOrGate, ACP.FeedForward.eval, ACP.FeedForward.toCircuit, ACP.nodeToCircuit_eval, BoolCircuit.Circuit.eval}
\difficulty{1}
Let $F$ be a Boolean feedforward circuit with input variables indexed by
$\mathrm{Fin}\,n$ and output nodes of type $\mathrm{out}$, equipped with a flag
assigning to each non-input node the label AND or OR and a witness that each gate has
finite arity. Suppose $F$ is an AND/OR circuit: every gate computes the conjunction of
the values on its input wires when its node is labelled AND, and their disjunction when
it is labelled OR. Then for every output node $o$ and every assignment $x :
\mathrm{Fin}\,n \to \bbf_2$, the tree-shaped Boolean circuit obtained by unrolling $o$
from the last layer down to the inputs, evaluated at $x$, agrees with the value that $F$
itself assigns to $o$ under the input assignment $x$.
\end{theorem}

\begin{theorem}[Size bound for the tree-unrolling of a feedforward circuit]
\lean{ACP.FeedForward.toCircuit_size_le}
\label{ACP.FeedForward.toCircuit_size_le}
\leanok
\uses{ACP.FeedForward, ACP.FeedForward.toCircuit, ACP.nodeToCircuit_size_le, BoolCircuit.Circuit.size, BoolCircuit.Circuit.toFeedForward}
\difficulty{1}
Let $F$ be a layered feedforward circuit over the Boolean alphabet with input type
$\mathrm{Fin}\,n$ and output type $\mathrm{out}$, of depth $d$, in which each gate is
designated an AND gate or an OR gate. Suppose every gate has finite input arity, and
that there is a bound $k \in \bbn$ such that every gate reads from at most $k$ input
wires. Then for each output node $o$, the tree-unrolled Boolean circuit obtained from
$o$ has size (its total number of AND/OR gates and literal leaves) at most $(k+1)^{d}$.
\end{theorem}

\begin{definition}[Embedding a Boolean circuit as a feedforward circuit]
\lean{BoolCircuit.Circuit.toFeedForward}
\label{BoolCircuit.Circuit.toFeedForward}
\leanok
\uses{ACP.FeedForward, ACP.FeedForward.GateOp.id, BoolCircuit.Circuit, BoolCircuit.Circuit.depth, BoolCircuit.Circuit.eval}
Turns a circuit $C$ on $n$ inputs into a feedforward circuit
$\texttt{FeedForward}\ \texttt{Bool}\ (\mathrm{Fin}\,n)\ \mathrm{Unit}$ of depth
$C.\texttt{depth} + 1$: layer $0$ carries the $n$ input variables and every higher
layer carries a single wire. The gate from layer $0$ to layer $1$ computes
$C.\texttt{eval}$ from all inputs at once, and all later gates are identity wires
that pass the single Boolean value upward.
\end{definition}

\begin{theorem}[Non-input nodes of the circuit embedding compute the circuit]
\lean{ACP.Circuit.toFeedForward_evalNode_const}
\label{ACP.Circuit.toFeedForward_evalNode_const}
\leanok
\uses{ACP.FeedForward.evalNode, BoolCircuit.Circuit, BoolCircuit.Circuit.depth, BoolCircuit.Circuit.eval, BoolCircuit.Circuit.toFeedForward}
\difficulty{4}
Let $C$ be a Boolean circuit on $n$ variables and let $x : \mathrm{Fin}\,n \to \bbf_2$
be an assignment of Boolean values to the inputs of $C$. The feedforward embedding of
$C$ is a layered circuit of depth $\mathrm{depth}(C) + 1$ whose layer $0$ carries the
$n$ input variables and each of whose higher layers carries a single node. Then for
every layer index $m$ with $1 \le m \le \mathrm{depth}(C) + 1$ and every node $v$ of
layer $m$, the value assigned to $v$ by the feedforward embedding on input $x$ equals
$\mathrm{eval}(C, x)$, the value of $C$ on $x$.
\end{theorem}

\begin{theorem}[Correctness of the feedforward embedding]
\lean{ACP.Circuit.toFeedForward_eval}
\label{ACP.Circuit.toFeedForward_eval}
\leanok
\uses{ACP.Circuit.toFeedForward_evalNode_const, ACP.FeedForward.eval₁, BoolCircuit.Circuit, BoolCircuit.Circuit.eval, BoolCircuit.Circuit.toFeedForward}
\difficulty{2}
Let $C$ be a Boolean circuit on $n$ variables, and let $x : \mathrm{Fin}\,n \to \bbf_2$
be an assignment of Boolean values to those variables. Form the feedforward embedding of
$C$, a layered feedforward circuit with a single output wire, and let its unique output
value under $x$ be the Boolean it computes. Then this value equals $C(x)$, the value
obtained by evaluating the circuit $C$ directly on the assignment $x$.
\end{theorem}

\begin{theorem}[Depth of the feedforward embedding]
\lean{ACP.Circuit.toFeedForward_depth}
\label{ACP.Circuit.toFeedForward_depth}
\leanok
\uses{BoolCircuit.Circuit, BoolCircuit.Circuit.depth, BoolCircuit.Circuit.toFeedForward}
\difficulty{1}
Let $C$ be a Boolean circuit on $n$ variables, with depth $\mathrm{depth}(C)$ equal to
the length of the longest root-to-leaf path in $C$. Embedding $C$ as a layered
feedforward circuit — in which layer $0$ carries the $n$ input variables and each higher
layer carries a single wire — produces a circuit whose depth equals $\mathrm{depth}(C) +
1$.
\end{theorem}

\begin{theorem}[Size bound for the feedforward embedding of a Boolean circuit]
\lean{ACP.Circuit.toFeedForward_size_le}
\label{ACP.Circuit.toFeedForward_size_le}
\leanok
\uses{ACP.FeedForward.size, BoolCircuit.Circuit, BoolCircuit.Circuit.depth, BoolCircuit.Circuit.one_le_size, BoolCircuit.Circuit.size, BoolCircuit.Circuit.toFeedForward}
\difficulty{3}
Let $C$ be a Boolean circuit on $n$ variables, and let $\mathrm{size}(C)$ and
$\mathrm{depth}(C)$ denote its total number of nodes and its depth, respectively.
Writing $\widehat{C}$ for the layered feedforward circuit obtained by embedding $C$, the
number of non-input nodes of $\widehat{C}$ satisfies
\[
\mathrm{size}(\widehat{C}) \;\le\; \mathrm{size}(C)\,\bigl(\mathrm{depth}(C) + 1\bigr).
\]
\end{theorem}

%%% added by the blueprint pipeline
%%%%%%%%%%%%%%%%%%%%%%%%%%%%%%%%%%%%%%%%%%%%%%%%%%%%%%%%%%%%%%%%%%%%%%%%%%%%%%%
\section{Additional declarations}

\begin{definition}[Gate operation]
\lean{ACP.GateOp}
\label{ACP.GateOp}
\leanok
A single operation in a feedforward circuit over an alphabet $\alpha$: it consists
of an arity type $\iota$ together with a function $(\iota \to \alpha) \to \alpha$
that computes the gate's output value from the values supplied on its input slots.
\end{definition}
